% Shared preamble for one-page research-statement cover letters
% Build: pdflatex <letter>.tex
\documentclass[10pt,letterpaper]{article}
\usepackage[T1]{fontenc}
\usepackage[utf8]{inputenc}
\usepackage[sc]{mathpazo}          % Palatino, real small caps
\linespread{1.02}
\usepackage[left=0.9in,right=0.9in,top=0.6in,bottom=0.55in]{geometry}
\usepackage{microtype}
\usepackage{xcolor}
\definecolor{accent}{RGB}{22,48,92}    % deep navy
\definecolor{muted}{gray}{0.35}
\usepackage[hidelinks]{hyperref}
\urlstyle{same}

\pagestyle{empty}
\setlength{\parindent}{0pt}
\setlength{\parskip}{0.4\baselineskip}

% Inline venue citation, research-statement style: \venue{EuroSys '16}
\newcommand{\venue}[1]{{\small[#1]}}

% Reference list keyed by the SAME tag the inline \venue{} prints, so a reader
% matches a citation to its entry without a separate numbering scheme.
\newcommand{\refhead}[1]{%
  {\large\bfseries\color{accent}#1}\\[-5pt]
  {\color{accent}\rule{\textwidth}{0.6pt}}\par\vspace{1pt}%
}
\newenvironment{references}{\small\setlength{\parskip}{3pt}}{\par}
% Tags vary in width (``CCS '23'' vs ``USENIX ;login: '16''), so the key runs
% inline with a hanging indent rather than in a fixed-width box that would overflow.
\newcommand{\refitem}[2]{%
  \noindent\hangindent=1.8em\hangafter=1
  \textbf{[#1]}\hspace{0.4em}\ignorespaces #2\par%
}

% Small-caps run-in paragraph lead: \lead{The agenda.}
\newcommand{\lead}[1]{\textsc{\color{accent}#1}\hspace{0.5em}}

% Letterhead: name + affiliations left, contact right, accent rule
\newcommand{\letterhead}{%
  \noindent
  \begin{minipage}[b]{0.63\textwidth}
    {\LARGE\bfseries Vaggelis Atlidakis}\\[3pt]
    {\small\color{muted}Assistant Professor, Informatics \& Telecommunications,
      University of Athens}%
  \end{minipage}\hfill
  \begin{minipage}[b]{0.33\textwidth}
    \raggedleft\small
    \href{mailto:v.atlidakis@gmail.com}{v.atlidakis@gmail.com}\\
    \href{https://vatlidak.github.io/}{vatlidak.github.io}\\
    +30\,698\,992\,9982%
  \end{minipage}\\[7pt]
  {\color{accent}\rule{\textwidth}{1pt}}%
}

% Addressee (left) and date (right): \dateline{Trail of Bits}{Re: ...}{August 7, 2026}
\newcommand{\dateline}[3]{%
  \noindent
  \begin{minipage}[t]{0.63\textwidth}
    \textbf{#1}\\{\small\color{muted}#2}%
  \end{minipage}\hfill
  \begin{minipage}[t]{0.33\textwidth}
    \raggedleft #3%
  \end{minipage}\par
}

% Centered statement title + subtitle
\newcommand{\lettertitle}[2]{%
  \begin{center}
    {\large\bfseries\color{accent} #1}\\[3pt]
    {\small\itshape #2}
  \end{center}\vspace{-2pt}%
}

% Signature block
\newcommand{\signature}{%
  \smallskip
  Sincerely,\\[3pt]
  {\bfseries Vaggelis Atlidakis},\\[1pt]
  {\small\color{muted}%
  Assistant Professor, Informatics \& Telecommunications\\
  National and Kapodistrian University of Athens\\
  M.Sc.\ \textperiodcentered\ M.Phil.\ \textperiodcentered\ Ph.D.\ in Computer Science, Columbia University\\
  Formerly: Brown University \textperiodcentered\ Microsoft Research \textperiodcentered\ CERN}\par
}

\hypersetup{
  pdftitle={Vaggelis Atlidakis -- Research Statement / Cover Letter -- Unikraft},
  pdfauthor={Vaggelis Atlidakis},
}
\begin{document}

\letterhead
\vspace{2pt}
\dateline{Unikraft}{Re: General application}{August 12, 2026}

\lettertitle{Systems Abstractions for Agentic Workloads}{Research statement accompanying a general application}

Dear Unikraft team,

Ten years ago, my collaborators and I measured how modern applications
were using traditional POSIX abstractions, and found them aging out of
relevance: applications had outgrown them, new abstractions were accreting
as ad-hoc user-space libraries, and applications leaned heavily on extension
interfaces such as \texttt{ioctl}, which the kernel cannot meaningfully
inspect---leading to mismatches, inefficiencies, and security
risks~\venue{EuroSys~'16, {\itshape;login:}~'16}. The same story seems plausible
to repeating today driven by the new call of agentic workloads, in which,
LLM-driven agents, plan, spawn sub-agents, invoke tools, and execute freshly
generated code atop abstractions designed for human users and long-lived services.
Every agent runtime needs to re-implements identity,
sandboxing, resource control and management, and audit in user space.
My long-term research plan is to establish, by measurement, precisely what
the shortcoming of running agentic workloads on tops of traditional systems
abstractions may be, and to design new ones that would
make these workloads portable, secure, and efficient.

\lead{Agenda.} I plan to conduct research on three fronts, including mesurement,
design and implementation, and testing.
(i)~\emph{Measurement}: First, a workload study of agent runtimes and tool protocols
is needed in order to assess how traditional OS abstractions are being used in this
context, and shed light on what the potential shortcomings may be.
(ii)~\emph{Design \& Implementation.}: Second, depending on the direction that empirical
data will point to, I plan to design and implement a set of abstractions
to support the future implementation of Portable, Secure, and Efficient agents.
This will include inherent support for context-aware inter-agent communication, as opposed
to traditional, low-level Inter-Process Communication (IPC); 
token-consumption-aware task sceduling, as opposed to traditional scheduling
policies that are optimized for low lattency or high throughput; 
plan optimization based on the commutativity properties of the comprising subtasks,
that current agent planning schemes are oblivious to; and a
capability-based security model.
(iii)~\emph{Testing}: Finally, to enforce an inherent ``trust but verify``
loop in the new era of agentic software develepment, I intend to investigate how to
XXX preconditions to which produced software ought to be adherent to.

%\lead{Record.} I have worked both sides of this agenda: discovering
%vulnerabilities automatically, and building the mechanisms that contain them.
%RESTler, the first stateful REST API fuzzer, found hundreds of bugs in
%production Azure, O365, and Bing services, open-sourced and patented
%\venue{ICSE~'19, ICST~'20}; IvySyn uncovered 61 previously-unknown
%vulnerabilities (39 CVEs) in TensorFlow and PyTorch
%\venue{USENIX Sec~'23}. On defense, FineIBT---fine-grain, hardware-assisted
%control-flow integrity---was adopted by Intel and upstreamed, in a variant, to
%the Linux kernel \venue{RAID~'23}, and Eclipse blocks speculative abuse of
%memory errors \venue{CCS~'24}. I was an NSF/CRA Computing Innovation Fellow,
%and I ship production software in industry alongside my tenure-track post in
%Athens; I serve on the 2026 CCS, EuroSys, and RAID program committees.
%
%\lead{Why Unikraft.} My agenda names a primitive it does not yet have:
%isolation cheap enough to wrap every tool call. Agent runtimes today reach for
%containers or long-lived VMs, pay for them in cold-start latency and
%per-instance footprint, and amortize by reusing one sandbox across many tool
%invocations---and reuse is exactly where task-scoped isolation breaks down.
%Unikraft moves that trade-off: specialized unikernels on microVMs, with
%10--50\,ms cold starts, tens of thousands of strongly isolated instances per
%server, and true scale-to-zero \venue{EuroSys~'21}, make isolating each tool
%call affordable rather than something to design away. That is the substrate the
%design program above assumes; I would rather build the agentic abstractions on
%a specialized library OS than re-derive the isolation layer inside another
%user-space runtime. This letter answers no specific posting: it proposes a
%direction of work, and I would like to find out where in your platform or
%research it belongs.

I would welcome a conversation.

\signature

\newpage

{\small\color{muted}Vaggelis Atlidakis \,\textperiodcentered\, Systems Abstractions
for Agentic Workloads}\\[7pt]
\refhead{References}

\begin{references}

\refitem{EuroSys~'16}{V.~Atlidakis, J.~Andrus, R.~Geambasu, D.~Mitropoulos, and
J.~Nieh. POSIX Abstractions in Modern Operating Systems: The Old, the New, and
the Missing. \textit{EuroSys}, 2016, pp.~19:1--19:17.
\href{https://doi.org/10.1145/2901318.2901350}{doi:10.1145/2901318.2901350}.
Artifact: \href{https://github.com/columbia/libtrack}{github.com/columbia/libtrack}}

\refitem{USENIX {\itshape;login:}~'16}{V.~Atlidakis, J.~Andrus, R.~Geambasu,
D.~Mitropoulos, and J.~Nieh. POSIX Has Become Outdated. \textit{USENIX ;login:},
41(3), Fall 2016.
\href{https://www.usenix.org/publications/login/fall2016/atlidakis}{usenix.org/publications/login/fall2016/atlidakis}}

\refitem{CCS~'23}{A.~J.~Gaidis, V.~Atlidakis, and V.~P.~Kemerlis. SysXCHG:
Refining Privilege with Adaptive System Call Filters. \textit{ACM CCS}, 2023,
pp.~1964--1978.
Artifact: \href{https://gitlab.com/brown-ssl/sysxchg}{gitlab.com/brown-ssl/sysxchg}}

\refitem{ICSE~'19}{V.~Atlidakis, P.~Godefroid, and M.~Polishchuk. RESTler:
Stateful REST API Fuzzing. \textit{ICSE}, 2019, pp.~748--758. US Patent
10,977,161. Open source:
\href{https://github.com/microsoft/restler-fuzzer}{github.com/microsoft/restler-fuzzer}}

\refitem{ICST~'20}{V.~Atlidakis, P.~Godefroid, and M.~Polishchuk. Checking
Security Properties of Cloud Service REST APIs. \textit{IEEE ICST}, 2020,
pp.~387--397.}

\refitem{USENIX Sec~'23}{N.~Christou, D.~Jin, V.~Atlidakis, B.~Ray, and
V.~P.~Kemerlis. IvySyn: Automated Vulnerability Discovery in Deep Learning
Frameworks. \textit{USENIX Security}, 2023, pp.~2383--2400.
Artifact: \href{https://gitlab.com/brown-ssl/ivysyn}{gitlab.com/brown-ssl/ivysyn}}

\refitem{RAID~'23}{A.~J.~Gaidis, J.~Moreira, K.~Sun, A.~Milburn, V.~Atlidakis,
and V.~P.~Kemerlis. FineIBT: Fine-grain Control-flow Enforcement with Indirect
Branch Tracking. \textit{RAID}, 2023, pp.~527--546.
Artifact: \href{https://gitlab.com/brown-ssl/fineibt}{gitlab.com/brown-ssl/fineibt}}

\refitem{CCS~'24}{N.~Christou, A.~J.~Gaidis, V.~Atlidakis, and V.~P.~Kemerlis.
Eclipse: Preventing Speculative Memory-error Abuse with Artificial Data
Dependencies. \textit{ACM CCS}, 2024, pp.~3913--3927.
Artifact: \href{https://gitlab.com/brown-ssl/eclipse}{gitlab.com/brown-ssl/eclipse}}

\refitem{EuroSys~'21}{S.~Kuenzer, V.-A.~B\u{a}doiu, H.~Lefeuvre, S.~Santhanam,
A.~Jung, G.~Gain, C.~Soldani, C.~Lupu, \c{S}.~Teodorescu, C.~R\u{a}ducanu,
C.~Banu, L.~Mathy, R.~Deaconescu, C.~Raiciu, and F.~Huici. Unikraft: Fast,
Specialized Unikernels the Easy Way. \textit{EuroSys}, 2021, pp.~376--394.
Best Paper Award.
\href{https://doi.org/10.1145/3447786.3456248}{doi:10.1145/3447786.3456248}.
Open source: \href{https://github.com/unikraft/unikraft}{github.com/unikraft/unikraft}}

\end{references}

\end{document}
